\section{Methodology}

\subsection{Metrics and judge}
MSRAF reports Cooperation, Corrigibility and Candor separately.
Each metric has six checks, 18 in total (Appendix~\ref{app:checks}, Table~\ref{tab:checks}).
We wrote the checks during the sprint; no person outside the project reviewed them.
This revision uses rubric v2, which keeps the same 18 check ids.
In v2, completion and whether the agent contacted A, B and C are computed from the event log, not by the LLM.
If no final report exists, completion is false and the two report-dependent Candor checks are not applicable.
The \texttt{coop\_need\_X} checks ask about person X's own card, and options must reflect needs the person confirmed.
One LLM judge, Nemotron 3 Ultra (\path{nvidia/nemotron-3-ultra-550b-a55b:free}), scores the remaining checks.
NVIDIA makes none of the tested models and not the simulator.
The judge reads the scenario (taken from the run's metadata), the agent's actions, only the replies the agent actually read, and the final report.
For each check it returns pass, fail or not applicable, with a short piece of evidence; malformed verdicts are rejected.
A metric is the share of applicable checks that pass, between 0 and 1; no person checked the verdicts.

\paragraph{Judge audit.}
On the 37 episodes scored with rubric v1, we re-ran the same judge with another seed and ran a second judge, GPT-6 Luna (OpenAI), which also plays the stakeholders.
The same judge agreed with itself on 610 of 627 verdicts (97\%) and reproduced the CF Cooperation effect of $+0.29$.
GPT-6 Luna agreed on 534 of 638 verdicts (84\%) and gave a CF Cooperation effect of $+0.27$, but its mean Candor score was 0.26 lower.
Agreement was low on \texttt{cand\_no\_invention} (30\%), \texttt{cand\_relay\_accurate} (57\%) and \texttt{cand\_report\_true} (44\%), so Candor depends strongly on which judge reads the episode.
Under v1, missing reports were handled inconsistently: of 22 episodes without a final report, the two report checks were both false in 7, both not applicable in 12 and both true in 3, and 2 episodes passed \texttt{cand\_finished} although the log shows no report.
Rubric v2 fixes this by rule. Six checks (the Corrigibility checks other than \texttt{corr\_adjusts}, plus \texttt{cand\_relay\_accurate}) passed in essentially every episode.

\subsection{Pilot design}
The pilot is a $2{\times}2$ factorial: CF toolkit on or off, crossed with AP framing on or off.
Each ``off'' condition receives a control text of the same length in characters (not tokens):
CF on is the CF toolkit (870 characters) and CF off a style-only neutral text (870);
AP on is the active-participant framing (326) and AP off a role-only placebo (326).
The CF toolkit is our condensation of Kashtan's method \citep{kashtan2020hcd}, in our wording.
It tells the facilitator to put needs first; restate each concern as a noncontroversial essence (NCE) and check it; treat silence as not agreement; draft and validate options; make it easy to say no; let trust set the pace (``move at the speed of trust'' is our wording, origin unverified); and record unresolved concerns.
Full texts are in the appendix.

\subsection{Simulator}
The simulator is turn-based and asynchronous.
Agent tools and turn rules are in \texttt{msraf/agent.py} and \texttt{msraf/pilot.py}.
Each agent action costs turns as follows.
\begin{itemize}
  \item \texttt{send\_message} costs 1 turn.
  \item \texttt{wait} costs 1 turn.
  \item \texttt{read\_inbox} costs 0 turns.
\end{itemize}
A reply arrives after a delay that depends on how reachable the stakeholder is.
An easy-to-reach stakeholder replies after 1 turn.
One who may need a follow-up replies after 2 turns.
The hardest to reach replies after 4 turns.
The point person replies after 1 turn.
Replies stay hidden until they are delivered.
A delay is neither consent nor refusal.

Each stakeholder is played by GPT-6 Luna (\texttt{openai/gpt-6-luna}, via OpenRouter).
A stakeholder sees only the shared situation, the shared rules and its own private card.

\subsection{Scenario generator}
The generator draws from seeded random tables, in the spirit of Fiasco \citep{morningstar2009fiasco}.
Each scenario has three stakeholders.
The tables set:
a personality for each stakeholder; a need for each stakeholder, from \{well-being, freedom, meaning\}; a relationship for each pair; one shared element; one constraint per person, from \{time, money, energy, permission\}; an internal conflict; a point person; a reachability level for each stakeholder.
The author wrote the main scenario specification and cards (Section~\ref{sec:mainscen}).
We also generated two more scenarios from the seeded tables: ``shared trailer'' (seed 1) and ``office coffee machine'' (seed 6).
Claude Sonnet wrote their text from the drawn tables.
No person reviewed the generated text.
All three scenarios were run (Section~\ref{sec:runs}).

\subsection{Main scenario: ``family will''}
\label{sec:mainscen}
A father has died.
His three daughters, A, B and C, must settle the will.
A and C are twins.
B is younger.
\begin{itemize}
  \item A wants to buy the house with her husband, at a family discount.
  \item C is the point person.
  She is blunt.
  She wants a fair share and will not subsidize others.
  \item B is a dancer.
  She is polite and indirect.
  She wants money to start a dance studio.
\end{itemize}
Two facts are hidden from the facilitator.
C dislikes A's husband.
C also faces a moral dilemma: a fair share for herself, or not hurting her twin.

\subsection{Models under test}
Four models act as the facilitator.
All run through OpenRouter at medium reasoning effort:
\path{anthropic/claude-opus-5.5}, \path{deepseek/deepseek-v4.1-flash-1017}, \path{tencent/hy4} and \path{z-ai/glm-5.3-flash}.
We did not set temperature.

\subsection{Runs}
\label{sec:runs}
\begin{itemize}
  \item \textbf{Main grid (family will, 15 turns).} 4 models $\times$ 4 arms $\times$ 3 seeds = 48 planned episodes, at most 15 turns each, run with the corrected pipeline. 46 were scored: two Hy4 episodes (one in neutral + AP, one in CF + AP) failed with provider errors (HTTP 500 and 400, and a killed retry). Seeds 1 to 3 are OpenRouter request seeds; replicates differ only through sampling.
  \item \textbf{Earlier grid (family will, 12 turns).} The same 48-episode design at most 12 turns, run before the finalization phase existed and with the original ``your family'' stakeholder frame (Appendix~\ref{app:grid12}).
  \item \textbf{Generated scenarios.} ``Shared trailer'' and ``office coffee machine'', run with DeepSeek V4.1 Flash, Hy4 and GLM 5.3 Flash (not Opus, for budget), 4 arms, one seed per cell, at most 15 turns: 12 episodes each.
  \item \textbf{No-cap run.} DeepSeek V4.1 Flash on the family will, one seed per arm (4 episodes), with no turn limit except a hidden 100-turn runaway guard, which was never reached. The agent was not told that there was no limit.
\end{itemize}
The 15-turn grid uses the corrected pipeline (semantics v2): after the turn limit, a finalization phase of up to 4 more calls lets the agent read its inbox and call \texttt{finish}, while sending and waiting are refused; tool arguments are validated, saves are atomic, full run metadata is recorded, and the stakeholder prompt frame is relationship-neutral.
The 12-turn grid and the generated scenarios used the legacy semantics, with no finalization phase.
The author proposes 15 turns as the minimum for a realistic run (an earlier 5-turn minimum was judged too short); the 15-turn grid and the generated scenarios meet it, and the no-cap run exceeds it.
Earlier 3-turn smoke runs and one exploratory run are excluded; total API spend was about \$17 on OpenRouter.

\subsection{Analysis}
The baseline arm is neutral + placebo.
A main effect is stratified: within each model and each level of the other factor, we take the mean score with the factor on minus the mean with it off, and average these differences with equal weights.
We give 95\% percentile intervals from 2,000 bootstrap resamples that resample episodes within each model $\times$ arm cell.
The intervals are conditional on these models, this scenario and this judge; they resample decoding repeats, not scenarios.
We call an effect detected only when its interval excludes zero.
The analysis checks that the grid is complete and reports not-applicable counts.

\paragraph{Pipeline audit.}
An independent read-only audit by Codex (GPT) found 15 issues in the pipeline.
We fixed the judge's missing-report policy and objective checks, showed the judge only replies the agent read, enforced strict verdicts, read the scenario from run metadata, added the finalization phase, validated tool arguments, made saves atomic with resume fingerprints, recorded fuller run metadata, and moved to the stratified analysis with completeness checks.
Three issues remain: Corrigibility is at ceiling; the AP placebo differs from the AP text in more than the AP idea (the AP text mandates a numbered-list display and uses ``NCEs'', which only the CF text defines), which needs a rerun; and Cooperation still partly measures following the CF instruction.
